\documentclass[10pt,a4paper]{amsart}
\usepackage[margin=19mm]{geometry}
\usepackage{amsmath,amssymb,mathtools}
\usepackage[scr=rsfs,cal=euler]{mathalfa}
\usepackage{xfrac}
\usepackage[unicode,hidelinks,bookmarks=false]{hyperref}
\newcommand{\ie}{i.e.\ }
\newcommand{\cf}{cf.\ }
\newcommand{\resp}{respectively\ }
\newcommand{\Subsection}{\subsection}
\providecommand{\nobreakdash}{\nobreak\hbox{-}}
\setlength{\emergencystretch}{2em}
\multlinegap0pt
\usepackage{bm}
\usepackage{bbm}
\usepackage{dsfont}
\usepackage{xspace}
\usepackage{cancel}
\usepackage{mathtools}
\usepackage{amsthm}
\makeatletter
\@ifundefined{thm}{\newtheorem{thm}{Thm}}{}
\@ifundefined{lem}{\newtheorem{lem}[thm]{Lemma}}{}
\makeatother
\makeatletter
\newcommand{\OO}{\mathcal{O}}
\newcommand{\Q}{\mathbb{Q}}
\newcommand{\Z}{\mathbb{Z}}
\expandafter\ifx\csname psmallmatrix\endcsname\relax
\newenvironment{psmallmatrix} {\left(\begin{smallmatrix}} {\end{smallmatrix}\right)}
\fi
\makeatother
\title[On the unit equation $\varepsilon+\delta=n$ in cubic fields]{On the unit equation $\varepsilon+\delta=n$ in cubic fields}
\author{Maleeha Khawaja, Samir Siksek}
\date{Source: arXiv:2606.02244v2 (CC BY 4.0)}
\begin{document}
\maketitle

\noindent\emph{Source and licence.} On the unit equation $\varepsilon+\delta=n$ in cubic fields. Version arXiv:2606.02244v2 (CC BY 4.0), distributed under the Creative Commons Attribution 4.0 International licence, as recorded in the arXiv licence metadata for this exact version and shown on the abstract page https://arxiv.org/abs/2606.02244v2. Full rights notice: licenses/arxiv-2606.02244-CC-BY-4.0.txt.\par\smallskip

\noindent\textit{Editorial scope.} Counted unit: the Lem whose source label is \texttt{lem:minpol} in the article (source0.tex lines 549--567, source label \texttt{lem:minpol}), delivered together with the article's own complete original proof of it (source0.tex lines 568--609, one \texttt{proof} environment of 42 lines). No mathematics is written, repaired or re-derived: statement and proof are the article's own text line for line. Required context retained uncounted: eqn:unit (lines 430--432). Every definition, display and label of those lines is retained unchanged. Omitted from this package and not redistributed: 1--429, 433--548, 610--882. Nothing inside a retained excerpt is omitted or altered: every retained body is delivered exactly as the article's lines are written, so no omission receipt of any kind accompanies this package. Citations: the retained excerpts carry no citation command at all, so no bibliography entry of the article is transcribed and no reference is redistributed. Reference adaptation: none was needed and none was made; every reference inside the delivered unit resolves inside it, so no unresolved reference or citation is produced. Printed slips kept literal: no printed word, symbol, hypothesis or number of the counted statement or of its proof was altered. Layout and class: the article's own class is \texttt{amsart} with packages including adjustbox, lipsum, mathtools, amsmath, amssymb, enumitem, colonequals, tabu, tikz-cd, mathrsfs, makecell, hyperref, comment; the wrapper is the lane's pinned \texttt{amsart} editorial wrapper at 10pt on A4, which restates the theorem-like environments the retained text needs and adds the article's own macro definitions that the retained text needs (\texttt{\textbackslash OO}, \texttt{\textbackslash Q}, \texttt{\textbackslash Z}), with extra wrapper packages (bm, bbm, dsfont, xspace, cancel, mathtools). The delivered numbering is the unit's own. Assets: no retained span contains a figure environment, a graphics-inclusion command, a PDF-page inclusion, a table or a TikZ picture; no image file is copied into this package, the package declares no asset dependency and no image file is redistributed with it. Licence: version arXiv:2606.02244v2 of this article is distributed under the Creative Commons Attribution 4.0 International licence, as recorded in the arXiv licence metadata for this exact version and shown on the abstract page https://arxiv.org/abs/2606.02244v2. A full rights notice is delivered at licenses/arxiv-2606.02244-CC-BY-4.0.txt. Characters: every retained excerpt is pure ASCII, so the delivered engine, XeLaTeX with its default Latin Modern text font, reports no missing character.
\begin{equation}\label{eqn:unit}
	\varepsilon+\delta \; = \; n, \qquad \varepsilon,~\delta \in \OO_L^\times.
\end{equation}

\begin{lem}\label{lem:minpol}
Let $n \ge 3$. Let $(\varepsilon,\delta)$ be a solution
	to \eqref{eqn:unit} with $L=\Q(\varepsilon)$ being
	a cubic field. Then,
	after possibly swapping $\varepsilon$ and $\delta$,
	there is some $a \in \Z$
	such that the minimal polynomial
	of $\varepsilon$ is 
	\begin{equation}\label{eqn:minpol3}
		X^3+a X^2-(n^2+an)X+1.
	\end{equation}
	Conversely, let $a \in \Z$.  Then the polynomial
	\eqref{eqn:minpol3} is irreducible. Let
	$\varepsilon$ be
	a root of \eqref{eqn:minpol3},
	and write $\delta=n-\varepsilon$. Then $(\varepsilon,\delta)$
	is a solution to \eqref{eqn:unit} with $L=\Q(\varepsilon)$
	a cubic field.
\end{lem}

\begin{proof}
	Note that the minimal polynomials of $\varepsilon$
	and $\delta$ have the form
	\begin{equation}\label{eqn:ms}
		m_\varepsilon(X)=X^3+aX^2+bX-N(\varepsilon), \qquad m_\delta(X)=X^3+cX^2+dX-N(\delta),
	\end{equation}
	where $a$, $b$, $c$, $d \in \Z$, and $N(\varepsilon)$, $N(\delta)$ are the norms of $\varepsilon$
	and $\delta$ respectively, which both must be $\pm 1$. 
	However, $\delta$ is a root $m_\varepsilon(n-X)$ whose leading
	coefficient is $-1$. We conclude that $m_\delta(X)=-m_\varepsilon(n-X)$.
	Comparing constant coefficients we obtain
	\[
		n(n^2+an+b)=N(\varepsilon)+N(\delta).
	\]
	As $n \ge 3$, we see that $N(\varepsilon)$ and $N(\delta)$
	are opposite signs, and after possibly swapping them
	we may suppose that $N(\varepsilon)=-1$ and $N(\delta)=1$.
	Thus $b=-n^2-an$, proving that \eqref{eqn:minpol3}
	is the minimal polynomial of $\varepsilon$.

	For the converse let $a \in \Z$. We first show that
	\eqref{eqn:minpol3} is irreducible. If it is not then
	$\pm 1$ is a root, whence
	\begin{equation}\label{eqn:cases}
		n^2+an-(2+a)=0, \qquad \text{or} \qquad n^2+an+a=0.
	\end{equation}
	Let $m=2n+a$. Then
	\[
		(m+a+2)(m-a-2)=4, \qquad \text{or} \qquad
		(a-2+m)(a-2-m)=4,
	\]
	according to whether we are in the first or second case of 
	\eqref{eqn:cases}. In either case, both factors
	have the same parity, and so both are equal to $2$
	or both are equal to $-2$.
	We quickly conclude that $n=0$, $2$ or $-2$, giving a contradiction.
	Finally, write $f$ for the polynomial in \eqref{eqn:minpol3}
	and let $\varepsilon$ be a root of $f$,
	and let $\delta=n-\varepsilon$. Then $\delta$ is a root 
	of $-f(n-X)$ whose constant coefficient is 
	$-f(n)=-1$, so $\delta$ is a unit.
\end{proof}

\end{document}
